\section{Introduction}
\label{sec:introduction}

A laboratory that has just expressed a new antigen and wants to know where its
surface-exposed B-cell epitopes are has good computational options, and a
practical problem in using them. The best current conformational predictors are built on
protein language models whose weights run to billions of parameters
\citep{lin2023esm}. Running one well needs either a workstation with a capable
accelerator or, far more commonly, uploading the sequence to a hosted service.
Neither is comfortable for every group, and accelerators are unevenly
distributed.
For antigens from unpublished constructs, from clinical isolates, or from an
industrial pipeline, the sequence itself is the sensitive object: sending it to
a third-party endpoint is a disclosure decision rather than a technical one.

The obvious middle path is to keep a small model on the local machine and call
the remote service only when the local model is unsure. Systems research has
established the idea under names such as split computing, model cascades and
early exiting \citep{matsubara2022split}. Its appeal in a screening setting is
easy to state: most antigens should be easy, the easy ones never leave the
device, and the hard minority gets the benefit of the larger model. What such a
scheme costs and returns is much less clear. A laboratory cares about three
quantities at once: how good the predictions are, how long screening takes, and
how much sequence leaves the building. We measure all three on the same runs
rather than arguing them separately.

We build a measurement harness that
runs a set of routing policies, from always-local through uncertainty-triggered
routing to always-remote, on identical inputs with identical seeds. It records
predictive quality, latency, data movement and an energy model per antigen.
Over \NumSeeds{} seeds on the BP3C50ID external split, the answer arrives with
intervals attached: routing on local uncertainty gives up nothing detectable in
quality against always-remote execution while cutting latency to roughly a
third and keeping most raw sequences on the device. It also gives an answer
nobody in this project expected. A routed antigen is charged twice: once for
the local pass that decides it is hard, once for the remote pass that then
processes it. Routing is therefore the most carbon-intensive policy in the
comparison. The green framing that usually accompanies edge inference does not
survive contact with the accounting, and Section~\ref{sec:routing-results}
gives the condition under which it would.

The paper also re-examines an earlier unpublished analysis of our own, which
reported that blending the scores of our local predictor with those of
BepiPred-3.0 \citep{clifford2022bepipred} improved on BepiPred-3.0 alone. The
mixing weight
and the decision threshold had both been chosen by scanning the evaluation
split and keeping the configuration that scored best.
\citet{kapoor2023leakage} identify that step as
the most common form of leakage in applied machine learning. It is easy to
commit, because a sensitivity analysis and a model-selection procedure look
superficially alike. We reproduce the original numbers under the original
protocol, to within \TunedBpThreeReplicationMaxDiff{}. We then re-estimate the
same quantity with weight and
threshold chosen inside the training folds of a leave-one-sequence-out loop.
The advantage does not survive on either dataset. On BP3C50ID the difference
against BepiPred-3.0 is no longer detectable and the MCC estimate changes
sign; on a larger IEDB-derived subset, which fuses against our own local
predictor rather than BepiPred-3.0, the difference is likewise indistinguishable
from zero. On both, the selection procedure places almost all of the weight on
the reference predictor rather than on ours. Neither dataset is large enough to
show that fusion cannot help, and we do not claim that it does not.

The literature this belongs to is accumulating fusion and ensemble results
faster than it scrutinises how their combination parameters were chosen, and
the corrective is cheap: the nested protocol here costs one extra loop over
sequences and needs no new data.

\paragraph{Contributions.} Concretely, this paper offers the following.
\begin{enumerate}
  \item A reproducible measurement harness for conformational B-cell epitope
  screening. It records predictive quality, latency, data movement and modelled
  energy together, per antigen and per seed, and states the provenance of each
  cost quantity (Section~\ref{sec:methods}).
  \item A \NumSeeds{}-seed characterisation of five local--edge routing
  policies on BP3C50ID, with bootstrap intervals and paired sign tests. It
  includes a closed-form condition locating where the known offloading energy
  threshold \citep{kumar2010cloud} falls for this task: routing increases
  modelled carbon once the routed fraction exceeds \CarbonBreakEvenRate{}, and
  a quality-tuned policy routes about twice that often
  (Section~\ref{sec:routing-results}).
  \item An instance, in conformational epitope prediction, of the
  selection-bias mechanism documented in bioinformatics since
  \citet{varma2006bias}: we reproduce a previously reported fusion advantage
  under its original, test-derived protocol, then re-estimate it leakage-free
  on two datasets, where it is no longer detectable and the estimate changes
  sign (Section~\ref{sec:fusion-results}).
  \item The observation that, given a free choice on either dataset, the
  selection procedure assigns our predictor almost no weight. We report the
  per-fold selections rather than summarising them away
  (Section~\ref{sec:weight-transfer}).
\end{enumerate}
